\section{Experiments}\label{sec:exp}

\paragraph{Protocol.} Every confirmatory experiment was pre-registered: the
criteria, thresholds and analysis scripts were frozen (with hashes) before
results were produced, and are not changed after the fact. Images are from
DIV2K validation (0801--0812) unless stated; $R_b$ is measured on luminance.
Exploratory analyses are labelled as such.

%-------------------------------------------------------------------
\subsection{The mismatch cost is predictable (E2)}\label{sec:exp-e2}

\paragraph{Setup.} 12 DIV2K images, five tasks (single image $\times2$,
$\times4$; sparse pixels 10\%, 25\%, 50\%). In each of the three bands we
multiply the true spectrum by $\kappa\in\{1/16,1/4,4,16\}$, reconstruct with a
GP (linear), and measure the lost recoverability in that band, giving 60
cells \src{E2§2}. Band energies are computed after a Hann window on
both signal and error; the window was not specified in the pre-registration
and was added after a one-image run showed boundary leakage (single-image
$\times4$ $b_3$ measured $R=0.56$ at 0.004\% transmittance). After the fix,
single-image $\times4$ $b_2$ measures $R=0.085$ against an a-priori 0.08
\src{E2§1}. Only the metric changed, not hypotheses or criteria.

\begin{table}[t]\centering\small
\caption{E2 pre-registered criteria (60 cells, 12-image means)
\src{E2§2; E4-T1}.}\label{tab:e2}
\begin{adjustbox}{max width=\linewidth}\begin{tabular}{lcc}\toprule
Criterion & Threshold & Result\\\midrule
a. Spearman(band-avg.\ prediction, measured) & $\ge0.8$ & \textbf{0.90}\\
b. $R<0.3$: over-est.\ $>$ under-est. & $\ge80\%$ & 4/4\\
b. $R>0.7$: under-est.\ $>$ over-est. & $\ge80\%$ & 6/7\\
c. median $|$pred.$-$meas.$|$ & $\le0.05$ & 0.038\\
d. Spearman(a-priori prediction, measured) & $\ge0.7$ & 0.79\\
\bottomrule\end{tabular}\end{adjustbox}
\end{table}

\begin{table}[t]\centering\small
\caption{Measured / band-averaged predicted cost $\Delta$ (fraction of band
energy) \src{E2§3}.}\label{tab:e2cells}
\setlength{\tabcolsep}{3pt}
\begin{adjustbox}{max width=\linewidth}\begin{tabular}{lccccc}\toprule
Task, band & $R$ & $\kappa{=}\tfrac1{16}$ & $\tfrac14$ & 4 & 16\\\midrule
SI$\times2$ $b_2$ & .80 & .230/.402 & .047/.101 & .067/.031 & .165/.056\\
SI$\times2$ $b_3$ & .11 & .084/.097 & .044/.058 & .194/.380 & \textbf{1.370}/2.648\\
SI$\times4$ $b_2$ & .085 & .060/.074 & .029/.045 & .150/.336 & \textbf{0.933}/2.590\\
Sparse 10\% $b_2$ & .37 & .197/.291 & .072/.144 & .110/.302 & .586/.976\\
Sparse 25\% $b_2$ & .70 & .226/.430 & .068/.139 & .094/.074 & .473/.149\\
Sparse 50\% $b_2$ & .90 & .148/.264 & .036/.046 & .024/.009 & .064/.015\\
\bottomrule\end{tabular}\end{adjustbox}
\end{table}

\paragraph{Results.} All four pre-registered criteria hold
(Table~\ref{tab:e2}). In low-$R$ bands under-estimation is nearly free
($\le0.06$) while a $16\times$ over-estimate costs 0.9--1.4 times the band
energy, i.e.\ worse than not reconstructing the band
(Table~\ref{tab:e2cells}). In high-$R$ bands the direction reverses. When
$R\approx0$ (single image $\times4$ $b_3$) neither direction costs anything,
because the observation carries nothing about the band \src{E2§3}.
The band-averaged prediction over-estimates large-$\kappa$ costs; for
decimation, evaluating Eq.~\eqref{eq:delta} per frequency with
Eq.~\eqref{eq:alias} and energy-weighting removes this Jensen gap (exploratory:
Spearman 0.99, median error 0.004; e.g.\ SI$\times4$ $b_2$, $\kappa=4$:
predicted 0.151, measured 0.150) \src{E2§4}.
The single counter-example to the asymmetry is sparse 10\% $b_1$ ($R=0.92$;
under 0.072 vs.\ over 0.110): random sampling spreads aliasing across bands,
outside the scalar model \src{E2§2}.

%-------------------------------------------------------------------
\subsection{Regime C: why ASISR decoders saturate (M2, E3)}\label{sec:exp-c}

\paragraph{Setup.} We compare four trained decoders that share the
EDSR-baseline encoder and the LIIF training protocol (DIV2K, 1000 epochs,
scale range $\times1$--$\times4$, seed 1): LIIF, LTE, SC-INR-eq and
SC-INR without the sinc footprint readout \TODO{confirm arm descriptions}
against two linear references, LMMSE with the generic spectrum and LMMSE with
the image's own (oracle) spectrum, on $256^2$ centre crops of the first 30
images of Set14, BSD100 and Urban100 \src{OUT7.1/h6c}.

\begin{table}[t]\centering\small
\caption{Single image $\times4$, $R_b$ measured on data
\src{OUT7.1/h6c}. Regime by $\Rlin$: $b_2$, $b_3$ are C. gen./orac.: LMMSE with generic/oracle spectrum; SC-eq/SC-ns: SC-INR-eq / SC-INR-nosinc.}
\label{tab:h6c}
\setlength{\tabcolsep}{3pt}
\begin{adjustbox}{max width=\linewidth}\begin{tabular}{llcccccc}\toprule
Set & band & gen. & orac. & LIIF & LTE & SC-eq & SC-ns\\\midrule
Set14 & $b_2$ & .16 & .25 & .56 & .57 & .57 & .56\\
 & $b_3$ & .09 & .09 & .20 & .20 & .20 & .18\\
BSD100 & $b_2$ & .13 & .22 & .38 & .38 & .38 & .38\\
 & $b_3$ & .06 & .06 & .13 & .13 & .13 & .11\\
Urban100 & $b_2$ & .11 & .27 & .50 & .50 & .51 & .51\\
 & $b_3$ & .02 & .03 & .14 & .14 & .15 & .13\\
\bottomrule\end{tabular}\end{adjustbox}
\end{table}

\paragraph{$\times4$.} In both regime-C bands the trained networks recover
far more than the oracle linear bound: 0.38--0.57 vs.\ 0.22--0.27 in $b_2$
(1.7--2.3$\times$) and 0.13--0.20 vs.\ 0.03--0.09 in $b_3$ (2--6$\times$)
(Table~\ref{tab:h6c}). This recovery cannot come from spectral adaptation of a
linear read-out (Prop.~\ref{prop:bound}); it comes from the learned prior.
The four decoders differ by at most 0.02 in either band.

\paragraph{Across scale factors (E3).} We repeat the measurement at
$\times2$, $\times3$, $\times8$ (27 cells: 3 scales $\times$ 3 bands $\times$
3 datasets) with the same networks; $\times8$ is out of the training range
\src{E3}. Pre-registered results (Table~\ref{tab:e3}): the a-priori
$\Rlin$ ranks the measured generic-LMMSE recovery with Spearman 0.952; in
10/12 regime-C cells all four networks exceed the oracle LMMSE and differ by
at most 0.03; in 9/9 regime-A cells the best network is within 0.05 of the
oracle LMMSE (largest gap $+0.016$); generic-LMMSE recovery decreases
monotonically with scale in 9/9 band--dataset pairs. The zones move as
predicted: $b_2$ is A at $\times2$ ($\Rlin=0.92$), B at $\times3$ (0.42) and
C at $\times8$ (0.011); $b_1$ enters B at $\times8$ (0.79)
\src{E3-key}. In all 12 C cells every network exceeds the oracle
(by $+0.002$ to $+0.244$) \src{E4-T2}.

\paragraph{Scale collapse (L1, pre-registered).} If regime C is set by the
observation rather than by the decoder, recovery should depend on frequency
relative to the LR Nyquist limit, $u=2s\nu$, not on absolute frequency. On
$240^2$ crops of the first 30 images per dataset, we measure $R(u)$ in seven
$u$ bins at $\times2/\times3/\times4$ for LIIF, LTE and SC-INR (seed 1) and
take the median range across scales, $D_{\mathrm{in}}$ \src{L1L2}. The
linear control collapses ($D_{\mathrm{in}}$ 0.019 / 0.006 / 0.011 on Set14 /
BSD100 / Urban100). Binning by $u$ instead of $\nu$ shrinks the range by an
order of magnitude in 9/9 model--dataset groups (0.013--0.061 vs.\
0.31--0.46). The main criterion, $D_{\mathrm{in}}\le0.05$ in $\ge8/9$ groups,
fails: 6/9, with all three models on Set14 at 0.055--0.061 (BSD100 and
Urban100: 0.013--0.019). This exceeds both the seed noise floor (median range
across seeds 0.003--0.004) and the linear control; the cause is not
diagnosed. We found no consistent deviation outside the training range
($\times6/8/12$; $D_{\mathrm{out}}>D_{\mathrm{in}}$ in only 6/9), and the
public models generalise only partly (GSASR and Thera pass, ContinuousSR
fails; 6/9) \src{L1L2}. Above the LR Nyquist limit, our three decoders
exceed the generic linear estimate by $+0.18$ to $+0.40$ and are almost
indistinguishable from each other \src{L1L2}.

\begin{table}[t]\centering\small
\caption{E3 pre-registered criteria \src{E3; E4-T1}.}\label{tab:e3}
\begin{adjustbox}{max width=\linewidth}\begin{tabular}{lcc}\toprule
Criterion & Threshold & Result\\\midrule
a. Spearman($\Rlin$, generic LMMSE) & $\ge0.8$ & 0.952\\
b. C: nets $>$ oracle, range $\le0.03$ & $\ge80\%$ cells & 10/12\\
c. A: best net $-$ oracle $\le0.05$ & $\ge80\%$ cells & 9/9\\
d. generic LMMSE monotone in scale & 9/9 & 9/9\\
e. (expl.) oracle$-$generic: B$>$A, B$>$C & -- & .078 vs .011/.040\\
\bottomrule\end{tabular}\end{adjustbox}
\end{table}

\paragraph{Reading.} In regime C the room for any decoder-side spectral
adaptation is set by what the learned prior already recovers, and the four
decoders sit within 0.03 of each other there. This is the band-level
counterpart of the 0.035\,dB saturation reported by \citet{nasir2026empirical}.
The ``within 0.03'' statement depends on the threshold: at 0.02 only 8/12 cells
pass \src{E4-E3sens}; see \S\ref{sec:limits}. Note also that regime A is
not where networks pull ahead either: there the oracle linear estimator is
already within 0.016 of the best network.

%-------------------------------------------------------------------
\subsection{Regime B: matching the spectrum decides (M3, M4)}\label{sec:exp-b}

\paragraph{Sparse pixels.} Sampling rates 5--50\%, DIV2K 0801--0812,
luminance; we compare linear estimation with the generic spectrum, with the
image's own spectrum, and DIP \citep{ulyanov2018dip} \src{OUT7.2/sp\_scan}. The true-spectrum linear estimator beats DIP at every rate
(26.39 vs.\ 25.19, 28.45 vs.\ 27.02, 32.66 vs.\ 30.84, 37.89 vs.\ 34.79\,dB at
5/10/25/50\%), and the generic spectrum trails it by 1.8--4.5\,dB
(Table~\ref{tab:sparse}). The space between the two linear estimators is the
whole headroom of spectrum estimation in this regime.

\begin{table}[t]\centering\small
\caption{Sparse-pixel reconstruction: $R_{b_2}$, $R_{b_3}$ and PSNR (dB) for
generic / true-spectrum linear / DIP \src{OUT7.2}.}\label{tab:sparse}
\setlength{\tabcolsep}{3pt}
\begin{adjustbox}{max width=\linewidth}\begin{tabular}{lccc}\toprule
Rate & $b_2$ & $b_3$ & PSNR\\\midrule
5\% & .17/.39/.16 & .10/.20/.05 & 24.59/26.39/25.19\\
10\% & .29/.56/.34 & .14/.31/.12 & 26.13/28.45/27.02\\
25\% & .59/.81/.69 & .29/.57/.38 & 29.27/32.66/30.84\\
50\% & .85/.95/.89 & .57/.82/.65 & 33.37/37.89/34.79\\
\bottomrule\end{tabular}\end{adjustbox}
\end{table}

\paragraph{1D EEG.} On 200 calibration segments of EEGMMIDB
\citep{goldberger2000physionet,schalk2004bci2000} with 10/25/50\% of samples
kept, the spectral prior helps most in $\theta$/$\alpha$ at 10--25\% and
almost nowhere at 50\%, where most bands are already A; e.g.\ $\alpha$ at
10\%: linear interpolation .22, population-average spectrum .30, true
spectrum .47 \src{OUT7.3}. \TODO{add H8 test-set (760 segments)
numbers from H8 result note; full table to appendix.}

%-------------------------------------------------------------------
\subsection{Regime A$\leftrightarrow$B: multi-frame SR (E1)}\label{sec:exp-ab}

\paragraph{Setup.} DIV2K 0801--0812, $\times4$, $T=16$ shifted frames
synthesised neutrally (footprints integrated on the $512^2$ original, ground
truth by $2\times2$ averaging), known registration, RGB, PSNR with an 8-pixel
border crop; $\sigma_n\in\{0,0.005,0.01,0.02\}$ \src{E1}. Arms: a
multi-resolution grid field (mrgrid); least squares with TV (ls\_tv); a
SuperF-style Fourier-feature MLP (ff\_$\sigma$); GPs with Gaussian spectra
(gp\_gauss\_$\sigma$), power-law spectra (gp\_power\_$\alpha$) and the
image's own spectrum (gp\_oracle); bilinear. Paired bootstrap 95\% CIs over
images.

\begin{table}[t]\centering\small
\caption{E1: 12-image mean PSNR (dB), $T=16$ \src{E1§2,
e1\_result\_v2.json}.}\label{tab:e1}
\begin{adjustbox}{max width=\linewidth}\begin{tabular}{lcccc}\toprule
Arm & $\sigma_n{=}0$ & .005 & .01 & .02\\\midrule
mrgrid & 31.39 & 28.36 & 22.57 & 14.36\\
ls\_tv1e-4 & 31.16 & 27.37 & 21.24 & 13.78\\
ls\_tv1e-3 & 29.98 & 29.61 & 28.64 & 24.84\\
ff\_3 & 27.94 & 27.92 & 27.88 & 27.74\\
ff\_10 & 28.86 & 28.79 & 28.61 & 28.01\\
gp\_gauss\_30 & \textbf{31.62} & 27.73 & 26.65 & 25.56\\
gp\_power\_2.0 & 31.19 & 29.37 & 27.37 & 26.10\\
gp\_power\_2.5 & 29.62 & 29.36 & 28.65 & 27.45\\
gp\_power\_3.0 & 28.27 & 28.21 & 28.04 & 27.50\\
gp\_oracle & 29.91 & \textbf{29.70} & \textbf{29.25} & \textbf{28.35}\\
bilinear & 24.65 & 24.64 & 24.59 & 24.41\\
\bottomrule\end{tabular}\end{adjustbox}
\end{table}

\paragraph{Results.} The direction predicted by Prop.~\ref{prop:mismatch}
holds: representations without spectral decay (mrgrid, weak TV) collapse
under noise, losing 17\,dB between $\sigma_n=0$ and $0.02$
(Table~\ref{tab:e1}), the over-estimation failure; the value of the true
spectrum over the best power-law GP grows with noise, $D$ from 0.33 to
0.85\,dB, difference $+0.51$ [0.13, 0.88] (P3 holds); and the winner moves
from a low-bias arm (gp\_gauss\_30) to gp\_oracle (P4 holds)
\src{E1§1}. At $\sigma_n=0$ a low-bias GP beats the best
Fourier-feature MLP by $+2.76$ [1.99, 3.56]\,dB.

Two pre-registered predictions fail \src{E1§1}. (P1) At
$\sigma_n=0.005$ low-bias representations no longer win: mrgrid $-$ ff\_10 $=
-0.43$ [$-1.01$, 0.14]\,dB; the switch happens earlier than predicted.
(P2) At $\sigma_n=0.02$ the generic power-law GP does not come within 0.6\,dB of
the oracle ($-0.85$ [$-1.10$, $-0.61$]). The winner margins at $\sigma_n=0$
and $0.005$ are not significant (+0.19 and +0.10\,dB, CIs include 0).

\paragraph{Why the switch came early.} The pre-registered prediction used one
zone label: at $T=16$, $R^{\mathrm{lin}}_{b_2}$ falls from 0.90 to 0.73 over
the noise range (Table~\ref{tab:map}), so $b_2$ is regime A at small noise and
low-bias arms were expected to win. It ignored each arm's mismatch $\kappa_b$
and ignored $b_3$. The D11 diagnostic below shows that both matter. Two further
notes. P4's switch at 0.005 partly reflects gp\_oracle's numerical noise floor
at $\sigma_n=0$ (it is optimal among linear arms for $\sigma_n>0$ by
Prop.~\ref{prop:bound}), so future switch predictions should compare practical
arms only \src{DER§2.2}. For the non-linear arms, mrgrid and ls\_tv1e-4 already
lose more at $\sigma_n=0.005$ than their noiseless lead over ff\_10, which is
the direct numerical reason P1 fails; their mechanism was not tested
\src{DER§2.2}.

\paragraph{Diagnostic on E1 (D11, pre-registered).} We stored per-frequency
error spectra of the GP arms on the E1 images and settings (CPU; the
non-linear arms mrgrid, ls\_tv and ff were not measured) \src{D11}. At
$\sigma_n=0.005$ the excess error of gp\_gauss\_30 over gp\_oracle splits
0.5\% / 44\% / 56\% over $b_1/b_2/b_3$. In $b_3$ the measured oracle
recovery is $R=0.22$ (regime B) with over-estimation $\kappa=4.5$; in $b_2$,
$R=0.80$ but $\kappa=16.4$ (Table~\ref{tab:d11}). The pre-registered rule
looked only at $b_2$; $b_3$ was already in regime B, and in $b_2$ a large
$\kappa$ is not free even at $R=0.8$, since the ``harmless over-estimation''
of regime A is the $R\to1$ limit. Pre-registered results: (a) the $b_3$ share
of gp\_gauss\_30's excess is 0.556 [0.465, 0.619], above 0.5 at the point
estimate but with a CI that includes it; (b) the band-wise predicted excess
$\sum_b E_b\,\Delta(R_b,\kappa_b)$ ranks the measured excess over 24 arm--noise
points with Spearman 0.856 [0.744, 0.918] (threshold 0.8, holds); (c) it does
\emph{not} pick the best non-oracle GP arm (1/3 noise levels, threshold 2/3):
the top arms are close and the band-averaged form carries a Jensen bias
(e.g.\ gp\_power\_2.0 at $\sigma_n=0.005$: predicted $1.40\times10^6$,
measured $3.9\times10^4$, unnormalised) \src{D11}. The concentration of excess
error on $R(\nu)<1/2$ frequencies (0.82 / 0.71 for the two wide arms vs.\ 0.18
for gp\_power\_3.0) has no discriminative value, because 78\% [74\%, 81\%] of
all frequencies satisfy $R(\nu)<1/2$ at this noise level; we do not use it
as evidence. So the band-wise cost $C(Q)$ (Eq.~\eqref{eq:crit}) ranks the GP arms but does
not select the winner; whether its a-priori version ranks equally well has not
been tested \src{DER§2.3}.

\begin{table}[t]\centering\small
\caption{D11: gp\_gauss\_30 vs.\ gp\_oracle, E1 multi-frame, $T=16$,
$\sigma_n=0.005$, 12-image mean \src{D11}.}\label{tab:d11}
\begin{adjustbox}{max width=\linewidth}\begin{tabular}{lccc}\toprule
Band & oracle $R_b$ & $\kappa_b$ & share of excess\\\midrule
$b_1$ & 1.00 & 1.02 & 0.5\%\\
$b_2$ & 0.80 & 16.4 & 44\%\\
$b_3$ & 0.22 & 4.5 & 56\%\\
\bottomrule\end{tabular}\end{adjustbox}
\end{table}

\noindent Non-linear arms (mrgrid, ls\_tv, ff) are covered only
qualitatively; $T=4/8$ (4 images) is descriptive only.


%-------------------------------------------------------------------
\subsection{Inside the network: SC-INR's footprint readout (F1)}\label{sec:exp-f1}

SC-INR's only component with a reliable benefit in our earlier ablations is
its footprint readout (analytic sinc integration over the target pixel); in
LTE the analogous input is the cell. F1 asks where in the band map this
component matters. We ablate it in two architecture pairs (LTE vs.\
LTE-NoCell; SC-INR vs.\ SC-INR-NoSinc) and measure the per-band loss
$L=R_{\mathrm{full}}-R_{\mathrm{ablate}}$ on 27 confirmatory cells
($\times2/\times3/\times4$ $\times$ 3 bands $\times$ 3 datasets), seed 1
\src{F1}. The pre-registered predictor is $g(R)=R(1-R)^2$, the
small-mismatch limit of Eq.~\eqref{eq:delta}.

\paragraph{Results (LTE pair).} $L$ is zero where $R\approx1$: mean $L$ over
regime-A cells is 0.0000, versus 0.0050 at $\times3$ $b_2$ (regime B)
\src{F1; E4-T1}. The loss concentrates in $b_3$ at $\times2$--$\times4$
and $b_2$ at $\times4$, where $\Rlin$ is only 0.01--0.10 (regime C) but the
network recovers 0.13--0.56. The a-priori predictor fails:
Spearman$(g(\Rlin),L)=0.411<0.5$ (standard Spearman with average ranks for
ties). Using the network's own recovery raises it to 0.490, still below the
threshold \src{E4-T1}.

\paragraph{Results (SC-INR pair).} With the matched pair (SC-INR vs.\
SC-INR-NoSinc, same encoder), the picture is the same: the A-regime criterion
holds, the a-priori predictor fails (Spearman$(g(\Rlin),L)=0.400$), the
network-recovery predictor also stays below threshold (0.492), and the
per-cell losses of the two architectures agree closely (Spearman 0.985)
\src{F1-corr2}. \TODO{A-regime mean $L$ and $\times3$ $b_2$ $L$ for the matched
SC pair: not in the corrected note.} All F1 numbers are seed 1 only; seeds 2/3
are evaluated under the G1 pre-registration below.


\paragraph{A closed-form band-order predictor (G1, pending).}
A training-free quantity, the relative gain difference between a one-pixel box
footprint and point sampling,
$\epsilon(\nu)=1-\operatorname{sinc}(\nu_x)\operatorname{sinc}(\nu_y)$,
grows with frequency and orders the bands. On seed 1, after the fact, its
band average $\epsilon_b^2$ ranks the losses better than $g(\Rlin)$ or
$g(\Rnet)$ (Spearman 0.664 for LTE) \src{G1 prereg §0}; this is post hoc and
not counted. A confirmatory test on seeds 2/3 (G1) was frozen beforehand.
We use $\epsilon_b^2$ only as a predictor of band order, not as a mechanism:
a linearised gain-mismatch model of the missing readout, with a factor $R$ or
$R(1-2R)^2$, was tested and rejected (pre-registered L2b: Spearman 0.62--0.65
vs.\ 0.82 for the direct measurement below), and across all closed-form
quantities we tried, multiplying by any $R$ factor gives no consistent gain
(0.60--0.71 vs.\ 0.65--0.68 without) \src{DER§3.4}.
\PH{G1 result on seeds 2/3 (both pairs, matched sc-inr / NoSinc).}

\paragraph{Predicting the loss from the full model (L2, pre-registered).}
One account is that the cell or footprint input tells the
decoder the scale factor, so a model without it behaves like the full model
averaged over the training scales. We test this without training: average
the full model's output over cells corresponding to $s'\in\{1.0,1.1,\dots,4.0\}$
and predict the ablation loss as $L_{\mathrm{pred}}=R_{\mathrm{full}}-R_{\mathrm{marg}}$
on the 27 F1 cells, seed 1 \src{L1L2}. The ranking matches:
Spearman$(L_{\mathrm{pred}},L)=0.822$ (LTE) and 0.824 (SC-INR), against
0.40--0.49 for the zone predictors of F1 and $+0.147/+0.149$ above the
footprint gain $\epsilon_b^2$ (0.676 on this evaluation) \src{L1L2}; with
log-uniform $s'$ it is 0.805/0.803. Two pre-registered criteria fail. The
magnitude matches only at $\times4$ and above: the median
$L_{\mathrm{pred}}/L$ is 0.59 (LTE) and 0.36 (SC-INR), outside $[0.5,2]$ for
SC-INR, because at $\times2/\times3$ the ablation loses more than
marginalisation predicts (e.g.\ BSD100 $\times3$ $b_3$: 0.038 vs.\ 0.008).
The ablated output is closer to the marginalised than to the full output in
73.9\% / 69.4\% of image--scale pairs (threshold 80\%), rising with scale
(LTE: 58\% at $\times2$, 99\% at $\times8$) \src{L1L2}. So the ablation loss can be predicted after the fact
from the full model's cell average, consistent with the readout supplying
scale information at large scale factors; at small scale factors the ablation
loses more than that, and no per-frequency linear account reproduces it
\src{DER§3.4}. This is correlational evidence; a
causal test needs a pure scalar-$s$ embedding arm, which requires training and
has not been run.

\paragraph{Reading.} The footprint readout pays off where the network has
recovered content, not where a linear estimator could: in the linear
regime C, the learned prior lifts these bands to intermediate recovery, and
only there does the readout matter. Of the explanations we tested, scale
marginalisation (L2) ranks the losses best; the zones supply only the
boundary condition (no loss at $\Rnet\approx0$ or $1$). The three regimes describe the inside of
an ASISR network only when assigned by $\Rnet$, which is post-hoc
(\S\ref{sec:zones}). We do not claim that $\Rlin$ predicts this component.

%-------------------------------------------------------------------
\subsection{Summary of predicted vs.\ observed winners (E4)}

\PH{Table of predicted vs.\ observed winners across all settings (E4-T2).
Ready to transcribe except the E1 rows (pending the corrected prediction) and
the F1 rows (seed 1 corrected; seeds 2/3 pending G1). Current tally from E4: one clear
mismatch, multi-frame $\sigma_n=0.005$.}
